\documentclass[12pt, a4paper]{article}
\usepackage[utf8]{inputenc}
\usepackage[T1]{fontenc}
\usepackage{amsmath, amssymb}
\usepackage{geometry}
\geometry{margin=2.5cm}
\usepackage{hyperref}
\usepackage{listings}
\usepackage{booktabs}
\usepackage{xeCJK}
\usepackage{tabularx}
\setCJKmainfont{PingFang TC}

\lstset{
  basicstyle=\ttfamily\small,
  breaklines=true,
  frame=single,
  backgroundcolor=\color{gray!10}
}
\usepackage{xcolor}

\title{用 Claude Code 建立研究生的自動化生活管理系統}
\author{Sheng-Yuan Chen}
\date{2026-04-14}

\begin{document}
\maketitle

\section*{前言}

身為研究生，同時要處理 RA 工作、TA 教學、自己的研究以及課程，資訊很容易分散在各個地方。這篇文章記錄我用 Claude Code 建立的一套自動化工作流，目標是：紀錄、整理、管理三位一體。

\section{系統架構}

整套工作流分為三層：

\begin{itemize}
  \item \textbf{每天}：Obsidian Daily Note（隨手記草稿）
  \item \textbf{每週}：\texttt{weekly\_notes/YYYY-WXX.md}（整理四個區塊）
  \item \textbf{週日}：自動上傳 Notion Weekly Review（完全不需動手）
\end{itemize}

\section{每週回顧系統}

\subsection{週記格式}

每週的筆記檔放在 \texttt{\textasciitilde/claude\_SY/weekly\_notes/YYYY-WXX.md}，固定四個區塊：
📚 Paper / 文獻、🔬 研究進度、🤖 AI Tools 學習、💡 想法 / TODO。

\subsection{自動上傳 Notion}

每週日 20:00，系統自動執行 \texttt{weekly\_review.py}：讀取週記 → 抓 git commits → 格式化 → Notion API 建立新頁面。

Cron 排程設定：

\begin{lstlisting}[language=bash]
0 20 * * 0 python3 ~/claude_SY/weekly_review.py \
  >> ~/claude_SY/weekly_review.log 2>&1
\end{lstlisting}

\section{Zotero 文獻管理}

\subsection{Better BibTeX 設定}

Citation key 格式：\texttt{[auth:lower][year]}，例如 \texttt{callaway2021}。

自動匯出設定：\texttt{File → Export Library → Better BibTeX → Keep updated → \textasciitilde/Desktop/NTU\_ECON/refs.bib}

LaTeX 引用：
\begin{lstlisting}[language=tex]
\bibliography{/Users/chenshengyuan/Desktop/NTU_ECON/refs.bib}
\end{lstlisting}

\subsection{Collections 結構}

\begin{itemize}
  \item 📥 Reading Queue（新 paper 先放這裡）
  \item 📐 Causal Inference（DiD \& Staggered、IV \& 2SLS、RDD \& Synth Control）
  \item 🤖 Causal ML（Double ML、Tree Methods、Neural Networks）
  \item 📊 Econometrics（Panel Data、Inference \& Testing）
  \item 🏫 Courses（MetricML、Advanced Econometrics）
  \item 🔬 RA\_jauer
\end{itemize}

\section{Obsidian Daily Note}

Vault 位置：\texttt{\textasciitilde/Desktop/NTU\_ECON/Notes/}，每天按 \textbf{Cmd+Shift+D} 開今天的 note。

Daily Note 模板涵蓋：今天讀了什麼、研究進度、AI Tools 心得、想法與問題、今日完成、明天待辦。

\section{Claude Code Skills 清單}

\begin{tabularx}{\textwidth}{lX}
\toprule
\textbf{指令} & \textbf{用途} \\
\midrule
\texttt{/new-research} & 建立新研究專案資料夾結構 \\
\texttt{/research-ideation} & 生成研究問題與實證策略 \\
\texttt{/interview-me} & 引導式問答正式化研究想法 \\
\texttt{/lit-review} & 系統性文獻搜尋 \\
\texttt{/review-paper} & 模擬 referee 論文評論 \\
\texttt{/new-post} & 發布學習筆記到網站 \\
\texttt{/data-analysis} & 完整 R 分析流程 \\
\texttt{/compile-latex} & 編譯 Beamer .tex \\
\texttt{/slide-excellence} & 簡報三合一審查 \\
\texttt{/commit} & git commit → PR → merge \\
\texttt{/learn} & 儲存本次技巧為 skill \\
\bottomrule
\end{tabularx}

\section{小結}

\begin{tabularx}{\textwidth}{lX}
\toprule
\textbf{工具} & \textbf{解決的問題} \\
\midrule
Obsidian Daily Note & 每天有地方記，不靠記憶 \\
weekly\_notes + cron & 每週彙整，自動上傳，不靠意志力 \\
Notion & 長期存檔，方便回顧 \\
Zotero + Better BibTeX & 文獻自動管理，.bib 永遠最新 \\
Claude Code Skills & 研究任務標準化，降低重複工作 \\
\bottomrule
\end{tabularx}

核心原則：\textbf{降低每一個環節的摩擦力}。越容易記錄，就越會記錄；越自動化，就越不依賴自律。

\end{document}
